\section{Introduzione}
Il progetto prevede la realizzazione di un’applicazione che permetta all’utente di scrivere semplici espressioni musicali testuali, ispirate al paradigma di Strudel/TidalCycles, e di trasformarle in sequenze temporali di eventi sonori. L’applicazione sarà pensata come un ambiente minimale per la composizione algoritmica: l’utente potrà definire pattern ritmici e melodici, modificarli durante l’esecuzione e osservare/ascoltare il risultato prodotto.

Il sistema non intende replicare integralmente Strudel, ma proporne una versione semplificata e realizzabile nel contesto del progetto PPS, concentrandosi soprattutto sulla modellazione funzionale dei pattern, sull’interpretazione della DSL e sulla gestione del flusso temporale degli eventi musicali.

